% Created 2026-07-06 Mon 11:49
% Intended LaTeX compiler: pdflatex
\documentclass{article}
               \usepackage{listings}
\usepackage{color}
\usepackage{amsmath}
\usepackage{array}
\usepackage[T1]{fontenc}
\usepackage{natbib}
               \usepackage{authblk}
\author{Hjalte Søberg Mikkelsen}
\affil{Steno Diabetes Center Copenhagen}
\RequirePackage{tcolorbox}
\definecolor{lightGray}{gray}{0.98}
\definecolor{medioGray}{gray}{0.83}
\definecolor{mygray}{rgb}{.95, 0.95, 0.95}
\newcommand{\mybox}[1]{\vspace{.5em}\begin{tcolorbox}[boxrule=0pt,colback=mygray] #1 \end{tcolorbox}}

\lstset{
keywordstyle=\color{blue},
commentstyle=\color{red},stringstyle=\color[rgb]{0,.5,0},
basicstyle=\ttfamily\small,
columns=fullflexible,
breaklines=true,
breakatwhitespace=false,
numbers=left,
numberstyle=\ttfamily\tiny\color{gray},
stepnumber=1,
numbersep=10pt,
backgroundcolor=\color{white},
tabsize=4,
literate={~}{$\sim$}{1}{ø}{{\o}}1{æ}{{\ae}}1{å}{{\aa}}1{Ø}{{\OE}}1{Æ}{{\AE}}1{Å}{{\AA}}1,keepspaces=true,
showspaces=false,
showstringspaces=false,
xleftmargin=.23in,
frame=single,
basewidth={0.5em,0.4em},
}

\usepackage[utf8]{inputenc}
\usepackage[T1]{fontenc}
\usepackage{graphicx}
\usepackage{longtable}
\usepackage{wrapfig}
\usepackage{rotating}
\usepackage[normalem]{ulem}
\usepackage{amsmath}
\usepackage{amssymb}
\usepackage{capt-of}
\usepackage{hyperref}
\date{}
\title{Constrained Risk Prediction}
\makeatletter
\newcommand{\citeprocitem}[2]{\hyper@linkstart{cite}{citeproc_bib_item_#1}#2\hyper@linkend}
\makeatother

\usepackage[notquote]{hanging}
\begin{document}

\maketitle
\newpage

\section*{Introduction}
\label{sec:org69a099e}

Accurate prediction of disease risk in the presence of competing risks remains a central challenge in clinical epidemiology. We target individualized predictions of the 5-year risk of type 2 diabetes (T2D) and simultaneously the 5-year risk of the competing risk of death. Common estimation methods, such as Fine-Gray regression and cause-specific Cox regression, target only one cause at a time and are not optimized for a given prediction time horizon.
We develop a loss-based estimation framework for the simultaneous prediction of multiple competing risks at a prespecified time horizon. Our approach provides a way to estimate a broad class of existing regression models for competing risks such that the resulting models are optimized for individualized risk prediction. Crucially, the framework enables seamless integration of flexible learner libraries and enables structural constraints during estimation. These may include coherence constraints that guarantee that the predicted risks respect the parameter space and domain-informed shape constraints (e.g., monotonicity in age). Estimation is achieved via constrained optimization of a differentiable loss function, such as the Brier score.


\section*{The Problem}
\label{sec:org0080255}

Assume we are in a competing risks scenario and want to predict the risks of experiencing each individual event at some timepoint \(\tau\). Denote \(T\) as the time between time0 and the date of an event and \(D\in\{0,1,...,K\}\) as the cause of the event, where \(D = 0\) is censoring. Let \(X = (X^1,...,X^p)\) be a p-dimensional vector of baseline covariates, which can both be numeric or categorical.\\[0pt]
Let \(\beta\) be a p-dimensional parameter vector and \(x = X^T\). Define the cumulative hazard for \(D = j\) as
\begin{align*}
\Lambda_j(t|x) = \exp(x\beta_j)\Lambda_{0j}(t)
\end{align*}
where \(\Lambda_{0j}\) is the baseline cumulative hazard for cause j, which is defined as \(\Lambda_{0j}(t) = \int_0^t \lambda_{0j}(s)ds\), where \(\lambda_{0j}\) is the baseline hazard, and \(t\in(0,\infty)\). Now define the survival function using the product integral as
\begin{align*}
S(t|x) = \prod_{s \leq t}\left(1-\sum_{i = 1}^K d\Lambda_i(s|x)\right),
\end{align*}
or using the exponential approximation
\begin{align*}
S(t|x) = \exp\left(-\sum_{i = 1}^K \Lambda_i(s|x)\right).
\end{align*}
We can then write the cumulative incidence function for cause \(j\), \(F_j(t|x) = P(T\leq t, D = j|X)\), as
\begin{align*}
F_j(t|x) = \int_0^t S(s-|x)d\Lambda_j(s|x).
\end{align*}
The problem I will look at is that current methods of estimation of the cause-specific risks estimates the cause-specific cumulative incidence functions seperately. The problem with this is that it can lead to wrong estimations, which can easily be verified as it can happen for \(\tau\) and some \(x^*\) that
\begin{align}
F_1(\tau|x^*)+F_2(\tau|x^*)>1.
\end{align}
I will in this project develope a method to estimate the risks simultaniously such that the above can't happen. MAYBE ALSO MENTION THE PROPORTIONAL HAZARDS ASSUMPTION, AND THAT BY USING BRIER-SCORE AS ESTIMATOR WE DON'T NEED THAT ASSUMPTION AS WE OPTIMIZE THE PARAMETERS FOR A GIVEN \(\tau\), AND DONT OPTIMIZE FOR THE WHOLE TIME PERIOD.


\section*{Model structure}
\label{sec:orge2b978e}
We are interrested in a fully parametric model that is optimized for one specfic timepoint. In Cox-models the baseline hazards are allowed to vary with time, but as we are only interrested in one timepoint, this is an unecessary complication. Therefore we choose to look at the parametric Cox-model with a constant baseline hazard, also known as the exponential Cox-model. In this model the baseline cumulative hazard function will be:
\begin{align*}
\Lambda_{0j}(\tau) = \int_0^{\tau} \lambda_{0j}ds = \tau\lambda_{0j},
\end{align*}
as \(\lambda_{0j}\) is now constant. The baseline hazard in this model can be viewed as the average baseline hazard up to \(\tau\).


\section*{Different loss functions}
\label{sec:org790d684}
To find the parameters that will make the model described in Model structure make the most accurate predictions, we need some loss function to minimize. We will here describe the usual maximum likelihood method and our new loss function based on the Brier-score. The loss functions will be described for general Cox-models and in depth for the exponential Cox-model.

\subsection*{Likelihood function}
\label{sec:org8b861bf}
Asume we have \(N\) datapoints and \(K\) different causes. The general parametric likelihood function in a competing risks setting can be written as:
\begin{align*}
L = \prod_{i = 1}^N\lambda_{j_i}(t_i|x_i)^{d_i}S(t_i|x_i),
\end{align*}
where \(d_i\) is 1 if an event was observed and 0 if \(i\) is censored, and \(j_i\) is which event occured for person \(i\). Using that the survival function is a product of the individual cause-specific survival functions, \(S(t_i|x_i) = \prod S_j(t_i|x_i)\), we can write the likelihood as:
\begin{align*}
L &= \prod_{i = 1}^N\lambda_{j_i}(t_i|x_i)^{d_i}\prod_{j = 1}^K\exp(-\Lambda_j(t_i|x_i))\\
&= \prod_{i = 1}^N\prod_{j = 1}^K\lambda_{j_i}(t_i|x_i)^{d_{ij}}\exp(-\Lambda_j(t_i|x_i)).
\end{align*}
This means that the likelihood function is a product of K likelihoods, one for each cause. We can then maximize this likelihood by maximizing each cause-specific likelihood function, as long as the parameters are different for each cause.

\subsubsection*{Likelihood estimating equations}
\label{sec:org93aef4c}
Assume two causes, only one covariate and constant baseline hazards (the model descriped in Model Structure). I then differentiate the log-likelihood functions with respect to the parameters. The negative log-likelihood function for cause \(j\) is then:
\begin{align*}
-\log(L_j) = \sum_{i = 1}^N-d_{ij}(\beta_jx_i+\log(\lambda_{0j}))+t_i\exp(\beta_jx_i)\lambda_{0j}
\end{align*}
and differentiated
\begin{align*}
\frac{d-\log(L_j)}{d\beta_j} = \sum_{i = 1}^N-d_{ij}x_i+t_ix_i\exp(\beta_jx_i)\lambda_{0j}
\end{align*}
\begin{align*}
\frac{d-\log(L_j)}{d\lambda_{0j}} = \sum_{i = 1}^N-\frac{d_{ij}}{\lambda_{0j}}+t_i\exp(\beta_jx_i)
\end{align*}


\subsection*{Brier score}
\label{sec:org8a984e9}
As our target in risk prediction is to find the model that minimizes the Brier-score, it makes sense to fit the models by choosing the parameters which minimizes the brier-score. By using the Brier-score as our loss function, we also optimize our parametres based on a single timepoint (the prediction horizon), which is of great interrest in risk prediction.
\\[0pt]
Assuming we have \(N\) datapoints of the form \((\tilde{T},\tilde{D},X)\) where \(\tilde{T}_i\in (0,\infty)\) is the observed time (time to first observed event), \(\tilde{D}\in\{1,2,...,K\}\) is the observed event and \(X\) is a p-dimensional covariate vector. The Brier-Score for predicting cause \(j\) at timepoint \(\tau\) can be written as
\begin{align*}
BS = \frac{1}{N}\sum_{i = 1}^N(F_j(\tau|x_i)-o_i(\tau,j))^2,
\end{align*}
where \(o_i(\tau,j)\) is the true outcome of person \(i\), i.e. \(o_i(\tau,j) = 1(\tilde{D}_i= j,\tilde{T}_i\leq \tau)\). Assume \(K\) different causes and assume no censoring (with censoring just add a 0 state such that we have \(K+1\) states and the setup will be similar to that of the joint survival super learner). We then want to find the parameters that minimize the Brier-score or all causes i.e.
\begin{align*}
BSS = &\arg\min_{\boldsymbol{\beta}\in\mathbb{R},\boldsymbol{\Lambda_0}\in\mathbb{R}^+}\sum_{j = 1}^K\frac{1}{N}\sum_{i = 1}^N(F_j(\tau|x_i,\boldsymbol{\beta},\boldsymbol{\Lambda_0})-o_i(\tau,j))^2\\
= &\arg\min_{\boldsymbol{\beta}\in\mathbb{R},\boldsymbol{\Lambda_0}\in\mathbb{R}^+}\frac{1}{N}\sum_{j = 1}^K\sum_{i = 1}^N\left(\int_0^{\tau} S(s-|x_i,\boldsymbol{\beta},\boldsymbol{\Lambda_0})d\Lambda_j(s|x_i,\boldsymbol{\beta}_j,\Lambda_{0j})-o_i(\tau,j)\right)^2,
\end{align*}
where \(\boldsymbol{\beta}\) are the beta coeffecients for all causes and \(\boldsymbol{\Lambda_0}\) are the baseline hazards, with j indicating the coeffecients for the specific cause. BSS = Brier-score-sum. This is actually just the Brier-score as it was written by Brier in the orginal paper. 



\subsubsection*{Find estimating equations in Parametric Cox with exponential baseline hazard}
\label{sec:org71aca03}
Assume we are in a competing risks setting with two causes and one covariate \(x\), and we are interrested in predicting the risk of experiencing the events before or at timepoint \(\tau\) using a cox model and let \(\phi = \{\beta_1,\beta_2,\Lambda_{01},\Lambda_{02}\}\). We will now look at the exponential Cox-model and are interested in solving this:
\begin{align*}
BSS=\arg\min_{\beta_1,\beta2\in\mathbb{R},\lambda_{01},\lambda_{02}\in\mathbb{R}^+}\frac{1}{N}\sum_{j = 1}^2\sum_{i = 1}^N\left(g_j(\phi)-o_i(\tau,j)\right)^2.
\end{align*}
First I rewrite \(g_1(\phi) = \int_0^{\tau} S(s-|x,\boldsymbol{\beta},\boldsymbol{\Lambda_0})\lambda_1(s|x,\boldsymbol{\beta}_1,\Lambda_{01})ds\) from the Brier-score written above (note that \(g_1\) is equal to \(F_1\) for a fixed timepoint \(\tau\)):
\begin{align*}
g_1(\phi)&= \int_0^{\tau}\exp\left[-\exp(\beta_1x)\Lambda_{01}(s-)-\exp(\beta_2x)\Lambda_{02}(s-)\right]\exp(\beta_1x)\lambda_{01}ds\\
&=\exp(\beta_1x)\lambda_{01}\int_0^{\tau}\exp\left[-\exp(\beta_1x)s^-\lambda_{01}-\exp(\beta_2x)s^-\lambda_{02}\right]ds\\
&=\exp(\beta_1x)\lambda_{01}\frac{\exp\left[\tau^-(-\exp(\beta_1x)\lambda_{01}-\exp(\beta_2x)\lambda_{02})\right]-1}{-\exp(\beta_1x)\lambda_{01}-\exp(\beta_2x)\lambda_{02}}\\
&=\exp(\beta_1x)\lambda_{01}\frac{1-\exp\left[\tau^-(-\exp(\beta_1x)\lambda_{01}-\exp(\beta_2x)\lambda_{02})\right]}{\exp(\beta_1x)\lambda_{01}+\exp(\beta_2x)\lambda_{02}}
\end{align*}
As we are in continous time \(\tau^-= \tau\). Then we differentiate \(g_1\) with respect to \(\beta_1\), \(\beta_2\), \(\lambda_{01}\) and \(\lambda_{02}\) (to get the result for \(g_2\) just swap all the 1's and 2's, also the ones on the left side of the equals sign):
\begin{align*}
\frac{\partial g_1(\phi)}{\partial\beta_1} = &\frac{x \lambda_{01} e^{x \beta_1}}{\left(\lambda_{01} e^{x \beta_1} + \lambda_{02} e^{x \beta_2}\right)^2}\\
&\cdot\left[
\lambda_{02} e^{x \beta_2}\left(1 - e^{-\tau\left(\lambda_{01} e^{x \beta_1} + \lambda_{02} e^{x \beta_2}\right)}\right)
+ \tau \lambda_{01} e^{x \beta_1}\left(\lambda_{01} e^{x \beta_1} + \lambda_{02} e^{x \beta_2}\right)e^{-\tau\left(\lambda_{01} e^{x \beta_1} + \lambda_{02} e^{x \beta_2}\right)}
\right]
\end{align*}
\begin{align*}
\frac{\partial g_1(\phi)}{\partial\beta_2} =\frac{xe^{\beta_1x}\lambda_{01}e^{\beta_2x}\lambda_{02}}{(e^{\beta_1x}\lambda_{01}+e^{\beta_2x}\lambda_{02})^2}\left[(1+\tau(e^{\beta_1x}\lambda_{01}+e^{\beta_2x}\lambda_{02}))e^{-(e^{\beta_1x}\lambda_{01}+e^{\beta_2x}\lambda_{02})\tau}-1\right]
\end{align*}
\begin{align*}
&\frac{\partial g_1(\phi)}{\partial \lambda_{01}}
=\frac{e^{\beta_1 x}}{\left(e^{\beta_1 x}\lambda_{01}+e^{\beta_2 x}\lambda_{02}\right)^2}\\
&\cdot\left[e^{\beta_2 x}\lambda_{02}\left(1-e^{-\tau\left(e^{\beta_1 x}\lambda_{01}+e^{\beta_2 x}\lambda_{02}\right)}\right)+\tau e^{\beta_1x}\lambda_{01}\left(e^{\beta_1 x}\lambda_{01}+e^{\beta_2 x}\lambda_{02}\right)e^{-\tau\left(e^{\beta_1 x}\lambda_{01}+e^{\beta_2 x}\lambda_{02}\right)}\right]
\end{align*}
\begin{align*}
\frac{\partial g_1(\phi)}{\partial\lambda_{02}} = \frac{e^{\beta_1x}\lambda_{01}e^{\beta_2x}}{(e^{\beta_1x}\lambda_{01}+e^{\beta_2x}\lambda_{02})^2}\left[(1+\tau(e^{\beta_1x}\lambda_{01}+e^{\beta_2x}\lambda_{02}))e^{-(e^{\beta_1x}\lambda_{01}+e^{\beta_2x}\lambda_{02})\tau}-1\right]
\end{align*}
Using this we can now differentiate the Brier-Score (BS) with respect to the betas to get the estimating equations:
\begin{align*}
\frac{\partial BSS(\phi)}{\partial\beta_1} &=\frac{1}{N}\sum_{j = 1}^K\sum_{i = 1}^N \frac{dg_j(\phi)}{\beta_1}2(g_j(\phi)-o_i(t,j))\\
&= 0
\end{align*}
\begin{align*}
\frac{\partial BSS(\phi)}{\partial\beta_2} &=\frac{1}{N}\sum_{j = 1}^K\sum_{i = 1}^N \frac{dg_j(\phi)}{\beta_2}2(g_j(\phi)-o_i(t,j))\\
&= 0
\end{align*}
\begin{align*}
\frac{\partial BSS(\phi)}{\partial\lambda_{01}} &=\frac{1}{N}\sum_{j = 1}^K\sum_{i = 1}^N \frac{dg_j(\phi)}{\lambda_{01}}2(g_j(\phi)-o_i(t,j))\\
&= 0
\end{align*}
\begin{align*}
\frac{\partial BSS(\phi)}{\partial\lambda_{02}} &=\frac{1}{N}\sum_{j = 1}^K\sum_{i = 1}^N \frac{dg_j(\phi)}{\lambda_{02}}2(g_j(\phi)-o_i(t,j))\\
&= 0
\end{align*}
Note that the \(\frac{1}{N}\) terms can be removed.


\section*{Estimation with constraint}
\label{sec:org8c7c6ca}

Assuming we have \(N\) datapoints, the Brier-Score for predicting cause \(j\) at timepoint \(\tau\) can be written as
\begin{align*}
BS_j = \frac{1}{N}\sum_{i = 1}^N(F_j(\tau|x_i)-o_i(\tau,j))^2,
\end{align*}
where \(o_i(\tau,j)\) is the true outcome of person \(i\) (1 if event j is the observed event and has happened at or before timepoint \(\tau\) and 0 otherwise). Furthermore assume only two competing risks (T2D and death). We then want to find the parameters that minimize this score i.e.
\begin{align*}
&\arg\min_{\boldsymbol{\beta}\in\mathbb{R}}\sum_{j = 0}^K\frac{1}{N}\sum_{i = 1}^N(F_j(\tau|x_i,\boldsymbol{\beta})-o_i(\tau,j))^2.
\end{align*}
In our setting \(K = 2\) where \(j = 0\) is the censored state and 1 and 2 are T2D and death. We need to minimize this score under some constraint, for example:
\begin{align*}
\sum_{j = 0}^K F_j(\tau|x,\boldsymbol{\beta})\leq 1 \text{ for all x},
\end{align*}
or
\begin{align*}
\frac{d}{d\beta_{age}}F_j(\tau|x_i,\boldsymbol{\beta})>0
\end{align*}
for some \(j\). In the T2D and death example \(j\in\{1,2\}\), as those risks should increase with age. In practise in may be easier to rewrite the minimization problem to:
\begin{align*}
&\arg\min_{\boldsymbol{\beta}\in\mathbb{R}}\sum_{j = 0}^K\frac{1}{N}\sum_{i = 1}^N(F_j(\tau|x_i,\boldsymbol{\beta})-o_i(\tau,j))^2+ I\left(\frac{d}{d\beta_{age}}F_j(\tau|x_i,\boldsymbol{\beta})\leq0\right),
\end{align*}
i.e. add the constraint to the loss function as a new term which will be 0 if the constraint holds and 1 otherwise which is large as we are using the Brier-score. We should probably multiply the indicator function with \(K+1\), as we sum \(K+1\) Brier-scores. \\[0pt]
If we want methods like Newton-Raphson to work we probably need to use a different method as above, as the loss-function need to be smooth. Say we have the constraint \(g(\beta)\leq 0\) then we can use the loss function:
\begin{align*}
&\arg\min_{\boldsymbol{\beta}\in\mathbb{R}}\sum_{j = 0}^K\frac{1}{N}\sum_{i = 1}^N(F_j(\tau|x_i,\boldsymbol{\beta})-o_i(\tau,j))^2+\lambda\max(0,g(\beta))
\end{align*}
where \(\lambda\) is some penalization parameter. 

\section*{Censoring}
\label{sec:orgc85fa83}
The BSS handles censoring well as is, and outperforms everything in the simulation study. Why? By this I mean we dont even add a censor state to the BSS.


\section*{Simulation plan}
\label{sec:org99826ee}

We want to show that using the Brier-score as the loss function, it can recover the true parameters of the data generating model. To do this we simulate data from an exponential Cox model, with 2 continuous and 1 binary covariate, where the effects of the covariates will be varied between different simulations. In each simulation the BS will be minimized for different time points (\(\tau\in\{0.5,1,2,4\}\)) and number of datapoints generated will vary (\(n\in\{100,200,500,1000\}\)). For each simulation a table will be produced with the columns: lossfunction, tau, n, true parameter, estimated parameter, mean difference between true par and est par, and standard deviation of the estimates. The rows will contain values for both the Brier-score as loss function and the likelihood. In all simulations the data will be right censored at \(t>5\) (Type 1 censoring), meaning if the first observed event happens after \(t = 5\), then \(\tilde{T}\) will be set to 5 and the event will be set to 0.

The above setup will be done both with and without different censoring processes (besides the type 1 censoring), where the censoring processes will be independent and dependent of the covariates.
\end{document}
